\section{Glossary}

This appendix gathers every term of the theory in one place, in plain words. It is meant to be read out of order, as a place to look something up rather than a chapter to work through. The structure mirrors the theory itself. First the two words that sit above everything, the substance and the whole. Then the eight atoms, the irreducible furniture of the base. Then the smaller words that live inside the atoms, the supporting vocabulary, the geometry and group terms, and finally the physics terms that emerge once the base is set in motion.

One discipline runs through all of it. The base is small and fixed. It is exactly eight atoms and nothing more. Every other word in this glossary either names one piece inside an atom, or names something that emerges from the eight once they run. The arrow of time, the attraction that binds matter, the self that persists, these are all emergent. They are not extra ingredients. They are consequences. Keeping that line sharp is the whole point of the lexicon, so the reader can always tell what is committed at the base from what is earned by the dynamics.

\subsection{The two words above the atoms}

Two words name the largest things in the theory. They are not atoms, because they are not parts. They are what the parts are made of and what the parts add up to.

\begin{definition}[Vibe]
A \emph{vibe} is the one substance, what everything is made of, and the name of the framework itself. Everything that exists is a vibe. The smallest vibe is one tone living on one site. Larger vibes, a dock, a self, the whole mesh, are vibes built out of smaller vibes. The word is scale-free. It names a bearer of experience at any size. Felt from inside, a vibe is the raw fact of something being experienced at all, a stir of feeling. A vibe is not a particle, and it is not the value it holds. The value is the tone. The bearer is the vibe.
\end{definition}

\begin{definition}[Flow]
The \emph{flow} is the whole mesh in motion, the eight atoms running as one. It is the universe taken as a totality unfolding, not a single step but the total process, every tone streaming through every dock across every beat. Where the knit is the law of one step, the flow is the living of that law forever. Written $F$ in the single-letter algebra, the capital that stands for the whole, distinct from the lowercase letters of the atoms.
\end{definition}

The opening claim of the theory can now be stated in these two words. The universe is a vibe mesh, and that mesh is a flow. Vibes vibing. Everything below unpacks what that sentence commits us to.

\subsection{The eight atoms}

The base is \textbf{eight atoms}. Each is a soft four-letter word, and each begins with a different letter, so each doubles as a clean single-letter symbol for the mathematics. Together they cover the universal furniture of any law. They say \emph{where}, they say \emph{what}, they say \emph{how the state changes} and \emph{how the space changes}, and they say \emph{when}. There are two laws rather than one because two different things evolve, the state living in the space and the space itself.

\begin{table}[ht]
\centering
\begin{tabular}{@{}llll@{}}
\toprule
atom & letter & plain meaning & answers \\
\midrule
\textbf{mesh} & m & the whole growing $\{3,4,3,4\}$ honeycomb, a weave of docks & where \\
\textbf{dock} & d & one here-now, the cell where tones berth for a beat & where \\
\textbf{site} & s & one of the $24$ directions out of a dock & where \\
\textbf{tone} & t & the substance value, ternary $-1 / 0 / +1$ & what \\
\textbf{lean} & l & the value order $-1 < 0 < +1$, charge and direction & what \\
\textbf{knit} & k & the state-law, collide then stream & how state changes \\
\textbf{wake} & w & the space-law, new docks born at the open edge & how space changes \\
\textbf{beat} & b & one discrete reversible tick of time & when \\
\bottomrule
\end{tabular}
\caption{The eight atoms of the base. Eight unique first letters, $m\,d\,s\,t\,l\,k\,w\,b$, assemble into $F$, the flow.}
\label{tab:eight-atoms}
\end{table}

\begin{principle}[The base is exactly eight]
The base is eight atoms and nothing else. Nothing else is committed. The arrow of time, the attraction that binds, and the self that persists are all earned from these eight by the dynamics. Adding any further ingredient to the base would be a cheat. If a phenomenon does not emerge from the eight, the honest report is the negative, not a ninth atom slipped in to force it.
\end{principle}

\subsubsection{mesh}

The \textbf{mesh} is space. It is the whole discrete space, the $\{3,4,3,4\}$ hyperbolic honeycomb, finite at any single beat and ever growing thereafter. It is not a static infinite stage laid down in advance. It is a weave of docks that grows at its own edge. The universe is a vibe mesh, so the mesh is the universe seen as the fabric every other thing lives in. In one sentence, the mesh is a weave of docks.

\subsubsection{dock}

A \textbf{dock} is a single place, a here-now, the geometric cell of the honeycomb. It is one place in the mesh. A dock has $24$ sites, the $24$ ways out of it, so a dock holds exactly $24$ tones, one per site. There is no rest slot and no twenty-fifth direction. A dock is a place of transit, not a store. Each beat, tones arrive at a dock as its fill, collide there by the knit, and depart as its will to flow on to neighbours. Felt from inside, a dock is a harbor of the present moment, a meeting then a moving on. What persists across time is never the dock but the self that threads through many docks. The band Dock is named for this atom, the atom of the here-now.

\subsubsection{site}

A \textbf{site} is one direction out of a dock, the directed slot where a single tone lives. The $24$ sites are the fixed set of all directions out of every dock, the same $24$ everywhere in the mesh. Concretely they are the $24$ roots of $D_4$, the vertices of the $24$-cell, the unit Hurwitz quaternions, arrayed as a burst of rays. Each site has an opposite among the $24$, the negated direction, so a tone that streams keeps a straight line of travel. A tone sits in one site and moves along it. There are $24$ sites per dock, never $25$.

\subsubsection{tone}

A \textbf{tone} is the value, the state a vibe is in. It is ternary, one of $-1$, $0$, $+1$, read as pain, peace, pleasure. The tone is the only attribute a vibe carries. A vibe has a tone, and its tone can change from beat to beat while the vibe itself persists. This is why vibe the substance and tone the value stay carefully apart. The vibe is the bearer, the tone is what it bears. Felt from inside, a tone is the flavor of a feeling, a push away at $-1$, a still rest at $0$, a draw toward at $+1$.

\begin{table}[ht]
\centering
\begin{tabular}{@{}lll@{}}
\toprule
tone value & reading & felt sense \\
\midrule
$-1$ & pain & a push away \\
$0$ & peace & a still rest \\
$+1$ & pleasure & a draw toward \\
\bottomrule
\end{tabular}
\caption{The single quale. One tone in $\{-1, 0, +1\}$. A whole dock of $24$ such tones is one hyper-color.}
\label{tab:one-tone}
\end{table}

\subsubsection{lean}

The \textbf{lean} is the value order, $-1 < 0 < +1$. The tones lean, a tilt from pain through peace to pleasure. The lean does two jobs at the base. It gives charge, the signed sum of a dock's tones. It gives the create move its direction, the way peace leans out a balanced pair $(+1, -1)$. The lean is the order through which the emergent arrow of time later runs, but the lean itself is just the ordering of the three values, committed at the base.

\subsubsection{knit}

The \textbf{knit} is the state-law, collide then stream. It is the local law that changes a dock's tones each beat, and it is the name of the law, never called by any other name. The knit has two phases. \emph{Collide}, the tones that meet head-on inside a dock interact by a fixed table. \emph{Stream}, each resulting tone then moves one dock along its site. The knit is reversible, charge-conserving, and local. It is not random and it is not global. To knit is to interlace, which is the collide, and to advance row by row, which is the beat-by-beat streaming, weaving the fabric of the mesh. The word pairs with mesh by design.

\subsubsection{wake}

The \textbf{wake} is the space-law, the growing edge. New docks are born at peace at the open, ever-expanding edge of the mesh, so the mesh grows. A wake is the expanding trail behind motion, and to wake is to come into being, so new docks wake into existence at the frontier. The openness of the edge is the bath. Radiation that streams outward falls behind the wake and never returns. In one sentence, the mesh grows at its wake, and the past is carried away.

\subsubsection{beat}

The \textbf{beat} is the unit of time, one tick. Each beat the knit runs once and the wake grows once. Time is the count of beats. There is no continuous clock underneath. The tick has no built-in direction, because the knit in the bulk is reversible. The direction of time, the arrow, is emergent, not committed at the base. Felt from inside, a beat is a heartbeat, the pulse of the world, a meeting then a moving on. A self persists across countless beats.

\subsection{The pillars and the single-letter algebra}

The eight atoms group into the universal furniture of a law. A law needs a where, a when, a what, and a how. Here the how splits in two, because two different things evolve, the state and the space.

\begin{table}[ht]
\centering
\begin{tabular}{@{}lll@{}}
\toprule
pillar & category & atoms \\
\midrule
where (space) & the stage & mesh, dock, site \\
when (time) & the tick & beat \\
what (substance) & the value & tone, lean \\
how the state changes & the state-law & knit \\
how the space changes & the space-law & wake \\
\bottomrule
\end{tabular}
\caption{The eight atoms grouped into the furniture of a law.}
\label{tab:pillars}
\end{table}

Because each atom takes a distinct first letter, the model reads as a clean symbol set. The flow is assembled from the eight.
\[
F \;=\; \langle\, \text{mesh},\ \text{dock},\ \text{site},\ \text{tone},\ \text{lean},\ \text{knit},\ \text{wake},\ \text{beat} \,\rangle
\;=\; \langle\, m,\ d,\ s,\ t,\ l,\ k,\ w,\ b \,\rangle.
\]
Here vibe names the substance the eight are made of, and $F$ names the whole they assemble into. Both sit above the atoms.

\subsection{The emergent words}

These are real and central terms, but they are not among the eight. They emerge from the eight once the flow runs. Keeping them out of the base is what keeps the theory honest.

\begin{definition}[Self]
A \emph{self} is a being, an "I". It is an emergent bound knot, a persistent, self-correcting, identity-bearing structure woven across many docks over many beats. It has three faces. \emph{Identity}, a winding that cannot be smoothly undone. \emph{Binding}, a knot held tight by the handedness of the sites. \emph{Agency}, it sheds disturbances as ripples into the open mesh, which is how it heals. A self is not a single dock and not a single tone. It is a pattern that keeps itself.
\end{definition}

\begin{definition}[The arrow]
The \emph{arrow} is the direction of time. It is emergent, not base. It arises from the lean, the value order $-1 < 0 < +1$, together with the create move and the wake acting as a low-entropy source at the open edge. The beat itself carries no direction. The arrow is what direction the beats acquire once the wake feeds the system from one side.
\end{definition}

\begin{definition}[The attraction]
The \emph{attraction} is the binding of matter, emergent and not base. It is binding by radiation, the shadow pressure of the active vacuum that holds a self together. It is the way the dynamics produce something that behaves like a force of attraction without any force being committed at the base.
\end{definition}

\subsection{The supporting words}

These smaller words name pieces inside the atoms or the everyday machinery of the dynamics.

\begin{table}[ht]
\centering
\begin{tabular}{@{}ll@{}}
\toprule
term & meaning \\
\midrule
\textbf{will} & a dock's outgoing sites, the projective push, the result of collide about to stream \\
\textbf{fill} & a dock's incoming sites, the receptive reading, perception, the tones just arrived \\
\textbf{rest} & the still state of a self, its localized knot of moving tones, the emergent rest mass \\
\textbf{base} & the foundation, the eight committed atoms, the discrete grain-level. Always "the base" \\
\textbf{code} & the concrete pattern, the knit's transition table, or a self's defining winding \\
\textbf{tree} & the branching hyperbolic bulk fanning out shell by shell, and selves nested in selves \\
\textbf{slot} & a code-level word only, the array position holding a site's tone \\
\bottomrule
\end{tabular}
\caption{The supporting vocabulary. Will and fill are the outgoing and incoming faces of one tone, not a second substance.}
\label{tab:supporting}
\end{table}

\subsubsection{The perception loop}

The beat, read experientially, is a perception loop. A dock \emph{fills}, receiving its neighbours' tones across its sites. It \emph{collides}, transforming them by the knit. It forms its \emph{will}, the projected result. The will \emph{flows} out to become the neighbours' fill. Then it repeats. All of it is still vibes. Fill and will are the incoming and outgoing faces of the one tone, not a separate kind of thing.

\subsection{Geometry and group terms}

These name the exact mathematical objects behind the geometry of the sites.

\begin{table}[ht]
\centering
\begin{tabular}{@{}ll@{}}
\toprule
term & plain meaning \\
\midrule
\textbf{line} & a direction and its opposite, the two head-on tones the collision acts on \\
\textbf{direction} & one of the $24$ $D_4$ root vectors $(\pm 1, \pm 1, 0, 0)$, a site \\
\textbf{opposite} & the negated direction, the reverse, needed for streaming \\
\textbf{degree} & the number of sites, $24$ \\
\textbf{charge} & the signed sum of tones, conserved by the knit \\
\textbf{$D_4$ roots} & the $24$ vectors $(\pm 1, \pm 1, 0, 0)$, the $24$ sites \\
\textbf{$24$-cell} & the regular $4$-polytope whose vertices are the $24$ directions \\
\textbf{Hurwitz quaternions} & the unit quaternions equal to the $24$ directions, a group of order $24$ \\
\textbf{$F_4$} & the sites' symmetry group, order $1152$, finite \\
\textbf{exceptional Jordan algebra} & the $27$-dimensional Albert algebra whose automorphisms are $F_4$ \\
\textbf{$600$-cell} & the finest regular $4$-polytope, $120$ vertices, smallest rotation about $36^\circ$ \\
\textbf{Hopf fibration} & the map from the $3$-sphere to the $2$-sphere, the soliton target \\
\bottomrule
\end{tabular}
\caption{Geometry and group terms behind the $24$ sites.}
\label{tab:geometry}
\end{table}

\subsection{The cosmology and substrate terms}

These name the larger structure of the growing mesh and its dynamics.

\begin{table}[ht]
\centering
\begin{tabular}{@{}ll@{}}
\toprule
term & plain meaning \\
\midrule
\textbf{pair table} & the committed collision, the create-flip-annihilate cycle plus hop plus inert \\
\textbf{create move} & peace $(0,0)$ brings out a balanced pair $(+1,-1)$, the active vacuum \\
\textbf{active vacuum} & empty space fluctuating with create-flip-annihilate pairs \\
\textbf{bath} & the open growing outside that swallows radiation for good \\
\textbf{cusp} & the point at infinity, its flat cross-section the physical $3$D space \\
\textbf{bulk} & the curved hyperbolic interior, a stronger bath, it disperses not binds \\
\textbf{horosphere} & the flat cross-section at the cusp, the $D_4$ lattice \\
\bottomrule
\end{tabular}
\caption{Cosmology and substrate terms. The arrow, the attraction, and the wake's role as a low-entropy source are emergent, covered above.}
\label{tab:cosmology}
\end{table}

\subsection{Physics terms, what emerges}

These name the physics that emerges once the flow runs and the grains are coarse-grained into smooth fields. None of them is committed at the base. Each is a consequence.

\begin{table}[ht]
\centering
\begin{tabular}{@{}ll@{}}
\toprule
term & plain meaning \\
\midrule
\textbf{Skyrmion} & a topological soliton, a knot in a direction field, charge equals its winding degree \\
\textbf{Hopfion} & a $3$D knotted or linked direction-field texture, the Hopf invariant \\
\textbf{instanton} & a $4$D winding texture \\
\textbf{Dzyaloshinskii-Moriya term} & a chiral twist that fixes a soliton's size, here the quaternion handedness \\
\textbf{Derrick obstruction} & pure exchange is scale-marginal, a soliton needs a stabilizer or it disperses \\
\textbf{soft mode, phonon} & a long-wavelength gapless collective excitation, the photon \\
\textbf{Poincare recurrence} & a closed reversible system returns near its start, so it cannot host a self \\
\textbf{symmetry restoration} & a finite group flowing to a continuous one in the coarse-grained infrared \\
\textbf{coarse-graining} & averaging many grains into a smooth field, where the real numbers come from \\
\textbf{triality, $8v\,8s\,8c$} & the split of the $24$ directions into a vector and two spinors \\
\textbf{$g$-factor} & the spinor magnetic moment, measured as $2$ from the Dirac structure of the sites \\
\textbf{Weyl fermion} & the massless limit, when the two chiralities decouple \\
\bottomrule
\end{tabular}
\caption{Physics terms. Each emerges from the eight atoms, none is part of the base.}
\label{tab:physics}
\end{table}

\begin{result}[The whole glossary in one line]
Everything in this appendix is either one of the eight atoms, a piece inside an atom, or something that emerges from the eight. The base is the eight, $m\,d\,s\,t\,l\,k\,w\,b$. The substance is the vibe and the whole is the flow. Everything else, the arrow, the attraction, the self, and all the physics, is earned, not assumed.
\end{result}
